\section{Details of graded unitarity constraints on $V_k(\sof_{2N})$ with $N$ odd} \label{app:VkSOodd}
In this appendix, we provide the details of the graded unitarity analysis for $V_k(\sof_{2N})$ with $N$ odd. 
\subsection{Low level constraints}
We begin with a bespoke analysis for conformal levels $L \leqslant 5$. To this end, note that the vacuum determinant in a fixed-root sector for $L = 1$ and $L = 2$ is given by,
\begin{equation} \label{eq: so2NoddL1to2}
\begin{split}
    &\text{det}~S_{\delta + \alpha} \ \propto \textcolor{blue}{(k - 0)}~, \\
    &\text{det}~S_{2\delta + \alpha} \propto \textcolor{blue}{(k - k_{N,1})}(k-k_{2N-2,1})^{\#} (\ldots \ldots)~. \\ 
\end{split}   
\end{equation}
Thus, up to $L = 2$, graded unitarity constraints require the level $k$ to be in the regime $-(2N-2) < k \leqslant -N+2$ or $k = 0$. The fixed-root vacuum determinant at $L = 3$ is given by,
\begin{equation} \label{eq: so2NoddL3}
\begin{split}
    &\text{det}~S_{3\delta + \alpha} \propto \textcolor{blue}{(k-k_{8,3})(k-k_{4,1})^3}(k-k_{5,1})^{\#}(\ldots\ldots) \qquad \text{when} \ \ N=5~, \\
    &\text{det}~S_{3\delta + \alpha} \propto \textcolor{blue}{(k - k_{2N-2,3})}(k-k_{N,1})^{\#_2} (\ldots \ldots) \qquad \qquad \ \text{when} \ \ N \geqslant 7~.\\ 
\end{split}   
\end{equation}
For $N = 5$, the primitive null states at $k = k_{8,3}$ and $k = k_{4,1}$ are components of the adjoint, and the $V_{\omega_4+\omega_5}$ representation. Hence, passing to a representation basis instead of a weight basis, graded unitarity constraints rule out $k \in (k_{8,3}, k_{5,1})$. For $N \geqslant 7$, graded unitarity constraints of \eqref{eq:grunitsignrule} rule out $k \in (k_{2N-2,3}, k_{N,1})$. 

At conformal level $L = 4$, the fixed-root vacuum determinant is given by,
\begin{equation} \label{eq: so2NoddL4}
\begin{split}
    &\text{det}~S_{4\delta + \alpha} \propto \textcolor{blue}{(k-k_{2,1})(k-k_{5,2})}(k-k_{8,3})^{\#}(\ldots\ldots) \qquad \text{when} \ \ N=5~, \\
    &\text{det}~S_{4\delta + \alpha} \propto \textcolor{blue}{(k - k_{N,2})}(k-k_{2N-2,3})^{\#_2} (\ldots \ldots) \qquad \qquad \text{when} \ \ N \geqslant 7~.\\ 
\end{split}   
\end{equation}
For $N = 5$, the primitive null states at $k = k_{2,1}$ and $k = k_{5,2}$ are components of the adjoint and the $V_{2\omega_1}$ representation, respectively. Determinant-based graded unitarity constraints of \eqref{eq:grunitsignrule} rule out $k \in (k_{2,1}, k_{5,2})$. Moreover, passing to a representation basis instead of a weight basis, graded unitarity constraints rule out the entire interval $k \in (k_{2,1}, k_{8,3})$. For $N \geqslant 7$, graded unitarity constraints of \eqref{eq:grunitsignrule} rule out $k \in (k_{N,2}, k_{2N-2,3})$. 

Finally, at $L = 5$, the fixed-root vacuum determinant is given by,
\begin{equation} \label{eq: so2NoddL5}
\begin{split}
    &\text{det}~S_{5\delta + \alpha} \propto \textcolor{blue}{(k - k_{8,5})}(k-k_{2,1})^{\#} (\ldots \ldots) \qquad \qquad \text{when} \ \ N=5~, \\    
    &\text{det}~S_{5\delta + \alpha} \propto \textcolor{blue}{(k - k_{2N-2,5})}(k-k_{N,2})^{\#} (\ldots \ldots) \qquad \text{when} \ \ N \geqslant 7.
\end{split}
\end{equation}
Thus, up to $L = 5$, graded unitarity constraints restrict $k$ to be in the regime,
\begin{equation} \label{eq:lowlevelsoOdd}
\begin{split}
    &-(2N-2) < k \leqslant -\frac{32}{5} \qquad \text{or} \qquad k \in \left\{-6, -\frac{16}{3}, -3, 0 \right\}\qquad \ \ \ \text{when} \ \ N=5~,  \\
    &-(2N-2) < k \leqslant k_{2N-2,5} \ \ \ \text{or} \ \ \ \  k \in \{k_{2N,4}, k_{2N-2,3}, k_{2N,2}, 0\}  \ \ \ \text{when} \ \ N \geqslant 7~. 
\end{split}
\end{equation}

\subsection{Structure of zeros and constraints on the level}
With the possible zeros of the level $L$ vacuum determinant for $N$ odd given by \eqref{eq: so2nprelimzero}, we can infer the precise structure of zeros for the level $L$ vacuum determinant. For $N \geqslant 9$ and $6 \leqslant L \leqslant N-2$ (the $N = 5,7$ cases in this regime are handled by the low level analysis above), the structure of zeros follows from \eqref{eq:soevenstr1} and is identical to that in Section \ref{sec:soevenstruczero}. As with the case of $V_k(\mathfrak{so}_{2N})$ with $N$ even, determining the structure of zeros requires additional care, owing to the qualitatively different behaviour for $L = (N-1,N)~\text{mod}~(2N-2)$ compared to other conformal levels. 

To this end, first note that for $N = 5$, the structure of zeros up to $L = N = 5$ has already been documented in the low level analysis, and graded unitarity constraints up to $L = 5$ are summarised in \eqref{eq:lowlevelsoOdd}. On the other hand, for $N \geqslant 7$, determinant-based constraints from graded unitarity up to $L = N-2$ restrict the level $k$ to be in the regime given by,
\begin{equation} \label{eq:constSOodduptoNm2}
\scalebox{0.93}{$\displaystyle
    -(2N-2) < k \leqslant k_{2N-2, N-2} \ \ \text{or} \ \ k \in \{k_{2N, N-3}, k_{2N-2, N-4}, \ldots, k_{2N-2, 3}, k_{2N, 2}, k_{2N-2,1}\}~.
$}
\end{equation}
Furthermore, at $L = N-1$ and $L = N$, the structure of zeros takes the following form.
\begin{itemize}
    \item At $L = N-1$, the most negative zero of the fixed-root vacuum determinant is given by $k = k_{2N-2, N-1} = k_{2,1}$, and this zero always occurs with multiplicity one. The second most negative zero is $k = k_{2N,N-1}$ and this zero also occurs with multiplicity one.\footnote{Note that from \eqref{eq:Lmultofk2N}, $k = k_{2N,N-1}$ first appears as a zero at $L = N-1$, and the corresponding primitive null state is a component of the $V_{2\omega_1}$ representation. Since the root multiplicity of $V_{2\omega_1}$ is one, the multiplicity of this zero in the fixed-root, conformal level $L=N-1$ vacuum determinant is one.} In addition, the third most negative zero is given by $k = k_{2N-2, N-2}$.  
    \item At $L = N$, the most negative zero of the fixed-root vacuum determinant is given by $k = k_{2N-2,N}$, and this zero always occurs with multiplicity one. Furthermore, the second most negative zero is given by $k = k_{2, 1}$.  
\end{itemize}
The novel difference between the analysis for $N$ even and $N$ odd, first appears at $L = N$, and can be traced back to the contrasting multiplicities in \eqref{eq:Lmultofk2Nm2}. The above structure of zeros, together with the low-level analysis in \eqref{eq:lowlevelsoOdd}, immediately implies that determinant-based graded unitarity constraints up to $L = N$ rule out all levels $k$ except,\footnote{At $L = N-1$, determinant-based constraints from \eqref{eq:grunitsignrule} only eliminate $k \in (k_{2,1}, k_{2N, N-1})$. However, passing to a representation basis rather than a weight basis, since the primitive null representations at $k=k_{2,1}$ and $k = k_{2N,N-1}$ are different, the entire region $k \in (k_{2,1}, k_{2N-2, N-2})$ is ruled out.}
\begin{equation} \label{eq:so4np2levellessNp1}
    -(2N-2) < k \leqslant k_{2N-2,N} \ \ \text{or} \ \ k \in \{k_{2,1},k_{2N-2,i}: i \in \mathcal{I}^{\rm odd}_N, k_{2N,j}: j \in \mathcal{J}^{\rm odd}_N \}~.
\end{equation}
where $\mathcal{I}^{\rm odd}_N$ and $\mathcal{J}^{\rm odd}_N$ are the sets,
\begin{equation}
\begin{split}
    &\mathcal{I}^{\rm odd}_N = \{1, 3, 5, \ldots, N-2\}~, \\ 
    &\mathcal{J}^{\rm odd}_N = \{2, 4, 6, \ldots, N-3\}~.
\end{split}
\end{equation}

For $N \geqslant 5$ and $L \geqslant N+1$, it can be shown by a direct computation that the fixed-root vacuum determinant takes the following form, \\
\\
\underline{\textbf{Case 1: } $L \neq (N-1,N)~\text{mod}~(2N-2)$}
\begin{equation}
\begin{split}
    &\text{det}~S_{L\delta+\alpha} \propto \textcolor{blue}{(k-k_{2N-2,L})^1}(k-k_{2N-2,L-2})^{\#}( \ldots \ldots) \qquad L \ \ \text{odd}~, \\
    &\text{det}~S_{L\delta+\alpha} \propto \textcolor{blue}{(k-k_{2N-2,L-1})^{\#}}( \ldots \ldots) \qquad \qquad \qquad \qquad L \ \ \text{even}~.
\end{split}
\end{equation}
\underline{\textbf{Case 2: } $L = (N-1,N)~\text{mod}~(2N-2)$}
\begin{equation}
\begin{split}
    &\text{det}~S_{L\delta+\alpha} \propto \textcolor{blue}{(k-k_{2N-2,L})^1}(k-k_{2N-2,L-1})^{\#}( \ldots \ldots)~.
\end{split}
\end{equation}
With this structure of zeros of the vacuum determinant, all-level determinant-based graded unitarity constraints are straightforward to implement. Note that for $L \neq (N-1,N)~\text{mod}~(2N-2)$, graded unitarity sign constraints of \eqref{eq:grunitsignrule} eliminate the region $k \in (k_{2N-2, L}, k_{2N-2, L-2})$ for $L$ odd, whereas, it does not give any constraints for $L$ even. Whereas for $L = (N-1,N)~\text{mod}~(2N-2)$, graded unitarity sign constraints rule out $k \in (k_{2N-2, L}, k_{2N-2, L-1})$. Combining this with the graded unitarity constraints up to $L = N$, \emph{i.e.}, \eqref{eq:so4np2levellessNp1}, we find that under the assumption of the standard $\mfR$-filtration \eqref{eq:AKMRfiltration} on $V_k(\mathfrak{so}_{2N})$ affine VOAs, fixed-root determinant-based graded unitarity rules out all levels $k$ except,
\begin{equation} \label{eq:VkSOoddgrunitresults}
k = \begin{cases}
        -(2N-2) + \frac{(2N-2)}{q}: (2N-2, q) = 1 \\
        -(2N-2) + \frac{(2N-2)}{q}: q \in 2\mathbb{N}+1 \ \ \text{and} \ \ (2N-2, q) \neq 1 \\
        -(2N-2) + \frac{N}{q} \qquad : q \in \mathbb{N} \ \ \text{and} \ \ 1 \leqslant q \leqslant \frac{N-3}{2} \\
        -(2N-2) + \frac{2}{q} \qquad \ : q \in 2\mathbb{N}+1 \\     
    \end{cases}
\end{equation}
Combining these graded unitarity restrictions on the level $k$ with those for $V_k(\mathfrak{so}_{2N})$ with $N$ even, we get the final result \eqref{eq:VkSOfullresults}.
